\chapter{The Ledgers}
% OUTLINE Ch9 "The Ledgers (back matter)". The citation ledger = the
% volume's spine (Kodo's second eyes committed); the physical-descriptions
% ledger; the face-grade ledger PER BEASTMASTER'S FLAGS: (flag 2) the two
% senses of "algebra" kept distinct; (flag 3) grades recorded for the faces
% THIS volume uses, no pre-emption of Vol 3's organization. Gate: Kodo
% (with second eyes on the citation ledger).

This book promised, on its first page, that nothing in it would ask the
reader to believe anything new, and a promise like that is audited, not
applauded. The audit is this chapter. Three ledgers: the names and dates
this book spent, chapter by chapter; the physical descriptions it borrowed
as openings; and the honest grades of the faces it used. If any sentence
of the preceding chapters cannot be traced to a row below or to the
reader's own bench practice, that sentence is a defect, and the book asks
to be held to exactly that standard.

\section*{The citation ledger}

\subsection*{Preface}
\begin{itemize}
\item Amontons, 1699 --- the laws of friction; Coulomb, 1785 --- the
separation of static from sliding. Back: the stuck block.
\item Lenard, 1902 --- the photoelectric observations; Einstein, 1905 ---
the quantum explanation. Back: the red lamp that never works.
\item Hooke, 1678; Young, 1807; Cauchy, 1822 --- the linear zone, its
modulus, and the stress--strain formalism. Back: the bar and its yield.
\item Joukowsky, 1898 --- the water-hammer surge. Back: the pipe that
bangs.
\item Classical surveying arithmetic, with satellite navigation as the
living instance. Back: the trilateration earn of the number four.
\end{itemize}

\subsection*{Chapter 1 --- The Local}
\begin{itemize}
\item Hausdorff, 1914 (\emph{Grundz\"uge der Mengenlehre}) --- the
neighborhood axioms. \item Kuratowski, 1922 --- the closure
axiomatization. \item Kelley, 1955 (\emph{General Topology}) --- the
assembled standard. \item Fr\'echet, 1906 --- abstract convergence before
the neighborhood axioms. \item Alexandroff and Hopf, 1935
(\emph{Topologie}) --- the separation taxonomy; Kolmogorov's T0
identification therein.
\end{itemize}

\subsection*{Chapter 2 --- The Chart}
\begin{itemize}
\item Riemann, 1854 --- the Habilitation lecture (published 1868 by
Dedekind). \item Whitney, 1936 --- the precise patchwork definition and
the embedding in dimension $2n{+}1$. \item Lee
(\emph{Introduction to Smooth Manifolds}) --- the modern standard.
\end{itemize}

\subsection*{Chapter 3 --- The Direction}
\begin{itemize}
\item Grassmann, 1844 (\emph{Ausdehnungslehre}) --- the vector space of
directed magnitudes. \item Ricci and Levi-Civita, 1901 --- the absolute
differential calculus; the two transformation laws as a calculus.
\item Dirac, 1930 (\emph{The Principles of Quantum Mechanics}) --- the
bra--ket pairing as the dual space in the physicist's house.
\item Lee (\emph{Introduction to Smooth Manifolds}) --- the tangent-space
formalization, carried over from the previous chapter's citation.
\end{itemize}

\subsection*{Chapter 4 --- The Machine}
\begin{itemize}
\item Misner, Thorne, and Wheeler, 1973 (\emph{Gravitation}) --- the
tensor as a machine with slots. \item Cauchy, 1822 --- the stress tensor,
the subject's founding physical instance. \item Cartan, 1899 --- the
exterior calculus; also the general boundary theorem. \item Green, 1828;
Gauss and Ostrogradsky, 1831; Kelvin, 1850, and Stokes, 1854 --- the
classical cases of the boundary theorem, with the attribution told whole
in the chapter.
\end{itemize}

\subsection*{Chapter 5 --- The Comparison}
\begin{itemize}
\item Christoffel, 1869 --- the symbols; the transformation of
differential expressions. \item Levi-Civita, 1917 --- parallel
displacement. \item Cartan, 1923 --- the connection freed from the
surrounding space. \item Koszul, via the standard texts --- the axiomatic
covariant derivative.
\end{itemize}

\subsection*{Chapter 6 --- The Loop That Fails to Close}
\begin{itemize}
\item Misner, Thorne, and Wheeler --- the loop pedagogy. \item Riemann,
1854 and 1861 --- the question, and the four-index apparatus of the Paris
essay (both posthumous). \item Ricci and Levi-Civita, 1901 --- curvature
as the commutator; the contractions bearing Ricci's name.
\end{itemize}

\subsection*{Chapter 7 --- The Ruler}
\begin{itemize}
\item Riemann, 1854 --- the line element. \item do Carmo
(\emph{Riemannian Geometry}); Lee (\emph{Riemannian Manifolds}) --- the
standard descent, including the fundamental theorem. \item Levi-Civita,
1917 --- the selected connection's name and picture. \item Minkowski,
1908 --- the signature of spacetime, one paragraph as promised.
\end{itemize}

\subsection*{Chapter 8 --- The Two Faces at the Boundary}
\begin{itemize}
\item Einstein, November 1915 --- the field equations, read on the
covariant face. \item Dirac, 1930 --- the quantum grammar, read on the
contravariant face. \item Madelung, 1927 --- the hydrodynamic form of the
Schr\"odinger equation; the owed face's nearest earned neighbor.
\item The founding text of the apparatus this series documents --- cited
as apparatus, once, for the program its opening pages set and fence.
\end{itemize}

\section*{The physical-descriptions ledger}

The openings of these chapters borrowed the reader's own practice, and
practice has provenance too: the locality of every bench report (the
working tradition of experimental physics); the observer's coordinates
(the same tradition, formalized by the geometers named above); velocity
as the pointed question (mechanics since Newton); the machine with slots
(Misner, Thorne, and Wheeler's pedagogy, borrowed with its name);
stress as force-per-area across an oriented cut (Cauchy); the carried
arrow and the loop experiment (the differential geometers' standard
picture, told in the gravitation textbook's form); the ruler and the
ball (the surveyor's and the analyst's common stock). None of these
descriptions is original here, and none carries a claim beyond the cited
mathematics it introduces.

\section*{The face-grade ledger}

This volume used, and here grades, exactly what it touched. Its two
\emph{inputs} were cited mathematics whole: topology (Chapter 1 --- the
neighborhood axioms through the T0 identification) and linear algebra at
a point (Chapter 3 --- spaces, duals, multilinearity). Its
\emph{subject} --- the thing constructed --- was differential geometry
(Chapters 2 and 4 through 7), assembled from those inputs with every
step cited. Its \emph{readings} were the two faces of Chapter 8: the
geometric description on the covariant face, the quantum on the
contravariant, each cited to its owner. And one face remains \emph{owed}
across the whole series: the fluid description, seeded in the founding
text, nearest-neighbored by Madelung, its dissipative body
undelivered --- named at its two natural sites in this book and spent at
neither.

One bookkeeping distinction, recorded so no reader conflates two senses
of one word: the linear algebra of Chapter 3 is a cited mathematical
input of this volume, and it is not the same thing as the arithmetic
this series' measuring apparatus uses internally, which is exact
rational arithmetic on counts and is documented in Volume 1. The two
live in different ledgers --- one belongs to the mathematics this book
teaches, the other to the instrument this series measures with --- and
this row exists to keep them apart.

\section*{The closing line}

The audit is the book. A volume built under the law that no original
idea may appear in it succeeds exactly when its back matter is
exhaustive and its chapters are traceable --- when the reader can put a
finger on any sentence and be pointed, by these ledgers, to a name, a
date, or their own laboratory habit. The mathematics belongs to the
names above. The assembly --- the order of chapters, the openings, the
insistence on prices paid in the open --- is all this book claims, and
an assembly is not an idea; it is a reading plan. The four soundings
belong to the volumes around this one. The ground now holds.
